\documentclass[11pt]{article}
\usepackage[margin=1in]{geometry}
\usepackage{parskip}
\usepackage[hidelinks]{hyperref}
\usepackage{enumitem}
\pagestyle{empty}
\begin{document}

\noindent Oren Livne\\
Pine Birch Analytics, LLC, 35 Kelinger Rd, Churchville, PA 18966-1033, USA\\
\href{mailto:oren.livne@gmail.com}{oren.livne@gmail.com} $\cdot$ ORCID 0000-0001-6700-483X

\vspace{1em}
\noindent 24 September 2026

\vspace{1em}
\noindent Dear Editor-in-Chief,

\vspace{0.5em}
Please consider the enclosed manuscript, \textbf{``Exact Solution of the Pawn Game
for up to Eight Pawns per Side,''} for publication in the \emph{ICGA Journal} as an
original research article.

The pawn game---chess played with pawns only, no kings---is a familiar teaching
position, but its game-theoretic value has not, to my knowledge, been settled in
print beyond the smallest cases. This paper solves it exactly for all
$n = 1,\dots,8$ pawns per side (left-justified for $n<8$) under both the game's
stated rule (no legal move $=$ loss) and the chess-stalemate variant (no legal move
$=$ draw), settling the previously open eight-pawn game: it is a first-player
(White) win under both rulesets.

I believe the work fits the ICGA Journal's scope and will interest its readership
for four reasons:

\begin{enumerate}[leftmargin=*]
\item \textbf{A completed game solution with a proved core.} Beyond the tabulated
values and optimal openings, the solver rests on a passed-pawn ``race'' evaluator
whose every verdict is a theorem (an exact result under the loss rule, and a
two-sided bound that drives alpha--beta cutoffs under the draw rule), together with
the game's symmetry group and a lock-free parallel transposition table.

\item \textbf{A clean, testable finding of independent interest:} en passant is
\emph{essential}, not incidental. Re-solving with en passant disabled flips four of
the sixteen values---every White win under the draw rule ($n = 4, 6, 8$) and the
seven-pawn game under the loss rule.

\item \textbf{A section written for chess players}, connecting the solution to
over-the-board themes (the unstoppable outside passer, the breakthrough sacrifice),
in the spirit of Philidor's maxim that the pawns are the soul of chess.

\item \textbf{A third rule variant that isolates the role of tempo.} Allowing a
player to pass (but not twice running) makes every game a first-player win under the
stated rule and a draw under chess stalemate, showing the pawn game's outcomes to be
governed by zugzwang. This variant was prompted by correspondence with FIDE Master
Carl Boor and is verified the same way as the base game.
\end{enumerate}

Correctness is established by fuzzing the move generator against an independent
engine and by cross-validating every search shortcut against unpruned search; as an
additional stress test, the tablebase never loses to Stockfish from a non-losing
side across 100 games per case. The solver, its independent reference oracle, all
result files, and the manuscript source are openly available in the accompanying
repository (\url{https://github.com/orenlivne/pawngame}), so every number in the
paper is reproducible.

The manuscript is original, has not been published previously, and is not under
consideration elsewhere. It has a single author. There are no conflicts of interest
and no funding to declare. I have no objection to single-blind review (the journal's
default) and can prepare an anonymized version should double-blind review be
preferred.

\vspace{0.5em}
Thank you for your consideration.

\vspace{1em}
\noindent Sincerely,

\vspace{0.5em}
\noindent Oren Livne\\
\href{mailto:oren.livne@gmail.com}{oren.livne@gmail.com}

\end{document}
